\section{自由幺半群的应用}

记号约定:$X$是集合.

\begin{lemma}{自由幺半群的自由积刻画}{}
	设$M=\bigast\limits_{x\in X}\N$,$\phi:X\to M$把$x$映到第$x$个分支中的$e_{x}$,则下列图表交换.
	\begin{center}
		\begin{tikzcd}
			X && M \\
			& {\Mo\left(X\right)}
			\arrow["{\phi_{i}}", from=1-1, to=1-3]
			\arrow["i"', hook, from=1-1, to=2-2]
			\arrow["{\exists!\phi}"', dashed, from=2-2, to=1-3]
		\end{tikzcd}
	\end{center}
\end{lemma}

\begin{proposition}{字的典范分解}{}
	设$\omega\in\Mo\left(X\right)$.
	\begin{enumerate}
		\item 存在唯一的$n\in\N$和$X\times\N$的序列$\left(\left(x_{i},m_{i}\right)\right)_{i=1}^{n}$使得$\forall 1\leq\alpha\leq n-1\left(x_{\alpha}\neq x_{\alpha+1}\right)$且$\omega=\prod\limits_{\alpha=1}^{n}x_{\alpha}^{m_{\alpha}}$.
		\item 若还有$\prod\limits_{\beta=1}^{m}\tilde{x}_{\beta}^{\tilde{m}_{\beta}}$,则$p\geq n$,并且$m=n$时$\forall 1\leq\alpha\leq n\left(x_{\alpha}=\tilde{x}_{\alpha}\wedge m_{\alpha}=\tilde{m}_{\alpha}\right)$.
	\end{enumerate}
\end{proposition}
